\subsection{INVARIANT HERMITIAN FORMS}

In this number, the letter k denotes the field R or C. Let V be a finite dimensional k-vector space, \(\Phi\) a separating \(^{2}\) positive hermitian form on V, G a group, g an R-Lie algebra, \(\rho: G \to \mathbf{GL}(V)\) a group homomorphism, \(\varphi: \mathfrak{g} \to \mathfrak{gl}(V)\) a homomorphism of R-Lie algebras.

\begin{enumerate}
\def\labelenumi{\alph{enumi})}
\tightlist
\item
  The form \(\Phi\) is invariant under G (resp. g) if and only if \(\rho(g)\) is unitary with respect to \(\Phi\) for all \(g \in G\) (resp. \(\varphi(x)\) is anti-hermitian \(^{3}\) with respect to \(\Phi\) for all \(x \in g\) ). Indeed, denote by \(a^{*}\) the adjoint of an endomorphism a of V with respect to \(\Phi\) ; for g in G, x in g, u and v in V, we have
\end{enumerate}

\[
\begin{array}{c} \Phi (\rho (g) u, \rho (g) v) = \Phi (\rho (g) ^ {*} \rho (g) u, v), \\ \Phi (\varphi (x) u, v) + \Phi (u, \varphi (x) v) = \Phi ((\varphi (x) + \varphi (x) ^ {*}). u, v); \end{array}
\]

thus, \(\Phi(\rho(g)u,\rho(g)v)=\Phi(u,v)\) for all \(u,v\) in V if and only if \(\rho(g)^{*}\rho(g)=\mathrm{Id}_{\mathrm{V}}\) ; similarly, \(\Phi(\varphi(x)u,v)+\Phi(u,\varphi(x)v)=0\) for all \(u,v\) in V if and only if \(\varphi(x)+\varphi(x)^{*}=0\) , hence the stated assertion.

\begin{enumerate}
\def\labelenumi{\alph{enumi})}
\setcounter{enumi}{1}
\item
  If the form \(\varPhi\) is invariant under G (resp. g), the orthogonal complement of a stable subspace of V is stable; in particular, the representation \(\rho\) (resp. \(\varphi\) ) is then semi-simple (cf.~Algebra, Chap. IX); moreover, for all \(g \in G\) (resp. \(x \in g\) ), the endomorphism \(\rho(g)\) (resp. \(\varphi(x)\) ) of V is then semi-simple, with eigenvalues of absolute value 1 (resp. with purely imaginary eigenvalues); indeed \(\rho(g)\) is unitary (resp. \(i\varphi(x)\) is hermitian, cf.~Algebra, Chap. IX).
\item
  Assume that \(k = \mathbf{R}\) . If \(G\) is a connected Lie group, \(\rho\) a morphism of Lie groups, \(\mathfrak{g}\) the Lie algebra of \(G\) and \(\varphi\) the homomorphism induced by \(\rho\) , then \(\varPhi\) is invariant under \(G\) if and only if it is invariant under \(\mathfrak{g}\) (Chap. III, §6, no. 5, Cor. 3).
\item
  There exists a separating positive hermitian form on V invariant under G if and only if the subgroup \(\rho(\mathrm{G})\) of \(\mathbf{GL}(\mathrm{V})\) is relatively compact (Integration, Chap. VII, §3, no. 1, Prop. 1).
\end{enumerate}

